%% Milestone Two — Results Tables
%% Table 1: Overall Comparative Performance Metrics
%% Table 2: Expansion Factor Stratified by Baseline Quality
%% Table 3: Synthesizability Sensitivity Analysis
%%
%% Filters applied: is_synth (depth=3), Tanimoto >= 0.3, Desirability >= 0.5
%% Diversity: 1 - mean pairwise ECFP4 Tanimoto (Morgan r=2, n=2048), sample n=300

% -----------------------------------------------------------------------
% Table 1: Overall Comparative Performance Metrics
% -----------------------------------------------------------------------
\begin{table}[ht]
\centering
\caption{Overall comparative performance metrics for each modification method
across 1,000 property specifications. Hit rate is the fraction of generated
variants passing both the synthesizability gate and property conservation
thresholds ($\tau_t \geq 0.30$, $\tau_d \geq 0.50$, AiZynthFinder depth 3).
Expansion factor is the mean number of passing variants per seed molecule.
Substructural conservation is the mean ECFP4 Tanimoto to the seed molecule.
Physicochemical conservation is the mean desirability geometric mean.
Drug-likeness is mean QED. Chemical diversity is $1 - $ mean pairwise
ECFP4 Tanimoto among hits.}
\label{tab:overall}
\begin{tabular}{llccccccc}
\toprule
Method & Description
  & \makecell{Hit Rate \\ (\%)}
  & \makecell{Expansion \\ Factor}
  & \makecell{Substructural \\ Conservation}
  & \makecell{Physicochemical \\ Conservation}
  & \makecell{Drug-likeness \\ (QED)}
  & \makecell{Chemical \\ Diversity} \\
\midrule
Baseline (Luo et al., 2025)        & Repeated PrexSyn sampling          & 3.5  & 1.50  & 0.656 & 0.780 & 0.623 & 0.875 \\
CReM (Polishchuk, 2020)            & Rule-based, fragment substitution   & 5.7  & 6.94  & 0.577 & 0.681 & 0.682 & 0.859 \\
mmpdb (Dalke et al., 2018)         & Rule-based, matched molecular pairs & 4.1  & 4.65  & 0.651 & 0.717 & 0.648 & 0.860 \\
JT-VAE (Jin et al., 2019)          & Learned, continuous latent space    & 5.1  & 5.89  & 0.499 & 0.836 & 0.706 & 0.846 \\
LibINVENT (Fialkova et al., 2022)  & Learned, scaffold decoration        & 0.3  & 1.09  & 0.559 & 0.858 & 0.636 & 0.771 \\
\bottomrule
\end{tabular}
\end{table}

% -----------------------------------------------------------------------
% Table 2: Expansion Factor Stratified by Baseline Quality
% -----------------------------------------------------------------------
\begin{table}[ht]
\centering
\caption{Expansion factor for each modification method, stratified by the
baseline synthesizability quality of the seed molecule. Tanimoto bins reflect
AiZynthFinder hit-rate quartiles of the seed molecules as generated by
PrexSyn. Baseline hit rate is the fraction of repeated PrexSyn samples that
pass the evaluation gates for seeds in that bin. Expansion factor is the
mean number of passing variants per seed molecule within the same bin
($\tau_t \geq 0.30$, $\tau_d \geq 0.50$, AiZynthFinder depth 3).}
\label{tab:expansion-stratified}
\begin{tabular}{lcccccc}
\toprule
Baseline Quality & \makecell{Tanimoto \\ Bin}
  & \makecell{Baseline \\ Hit Rate (\%)}
  & \makecell{CReM \\ Expansion}
  & \makecell{mmpdb \\ Expansion}
  & \makecell{JT-VAE \\ Expansion}
  & \makecell{LibINVENT \\ Expansion} \\
\midrule
Low Quality         & $< 0.50$        & 0.1 &  0.88 &  1.52 &  1.68 & 0.17 \\
Medium-Low Quality  & $0.50$--$0.70$  & 0.5 &  2.60 &  1.27 &  0.59 & 3.04 \\
Medium-High Quality & $0.70$--$0.85$  & 1.2 &  7.92 &  5.54 &  6.39 & 0.24 \\
High Quality        & $0.85$--$1.00$  & 1.7 & 16.91 &  9.83 & 14.12 & 0.96 \\
\bottomrule
\end{tabular}
\end{table}

% -----------------------------------------------------------------------
% Table 3: Synthesizability Sensitivity Analysis
% -----------------------------------------------------------------------
\begin{table}[ht]
\centering
\caption{Synthesizability sensitivity analysis across AiZynthFinder search
depths. Depth 1, 3, and 6 refer to the maximum number of retrosynthetic steps
permitted during the Monte Carlo tree search. Hit rate is the fraction of
variants passing all evaluation gates at each depth ($\tau_t \geq 0.30$,
$\tau_d \geq 0.50$). Absolute change is the percentage-point difference in
hit rate between depth 1 and depth 3. Relative change is that difference
divided by the depth-1 hit rate.}
\label{tab:synthesizability-sensitivity}
\begin{tabular}{lccccc}
\toprule
Method
  & \makecell{Depth 1 \\ Hit Rate (\%)}
  & \makecell{Depth 3 \\ Hit Rate (\%)}
  & \makecell{Depth 6 \\ Hit Rate (\%)}
  & \makecell{Absolute \\ Change (1--3)}
  & \makecell{Relative \\ Change (1--3)} \\
\midrule
Baseline (Luo et al., 2025)        & 1.9 & 3.5 & 3.9 & $+1.6$ pp & $+85\%$  \\
CReM (Polishchuk, 2020)            & 2.2 & 5.7 & 6.4 & $+3.6$ pp & $+166\%$ \\
mmpdb (Dalke et al., 2018)         & 1.5 & 4.1 & 4.7 & $+2.6$ pp & $+179\%$ \\
JT-VAE (Jin et al., 2019)          & 1.7 & 5.1 & 5.7 & $+3.4$ pp & $+198\%$ \\
LibINVENT (Fialkova et al., 2022)  & 0.3 & 0.3 & 0.4 & $+0.1$ pp & $+31\%$  \\
\bottomrule
\end{tabular}
\end{table}
